\section{Related Work}
\label{sec:related_work}

\subsection{Reflection and experiential memory}
\label{sec:related_reflection}

Language-agent memory first developed around verbal feedback and reusable experience. Reflexion stores natural-language reflections from prior attempts and reintroduces them as episodic memory, while Agent-Pro translates interaction failures into policy-level reflection and optimization \cite{shinn2023reflexion,zhang2024agentpro}. These methods show that an agent can improve without updating the base model, but their memory objects primarily record retrospective lessons or policy adjustments.

Cross-task systems move from individual reflection toward abstraction over multiple trajectories. ExpeL distills general insights from successes and failures and combines them with retrieval of successful experiences \cite{zhao2024expel}. Investigate-Consolidate-Exploit similarly consolidates cross-task experience before retrieving it for new tasks \cite{qian2024ice}. Together, these systems establish trajectory-to-text experience compression. They do not, however, require one frozen artifact to encode applicability, ordered execution, and rejection conditions for a debugging family.

\subsection{Guidelines, workflows, and procedural skills}
\label{sec:related_procedural}

Later work gives experience a more operational structure. AutoGuide induces context-aware guidelines from interaction trajectories and selects them according to the current state \cite{fu2024autoguide}. AutoManual maintains an instruction manual through rule addition, revision, merging, deletion, and validation \cite{chen2024automanual}. Agent Workflow Memory is the closest natural-language workflow precedent: it induces reusable workflows from agent trajectories in offline and online settings and retrieves them for later tasks \cite{wang2025workflowmemory}. These systems demonstrate that natural-language memory can express more than retrospective reflection, including conditional guidance and multi-step workflows.

A parallel line of work stores executable or structured skills. Voyager builds an executable skill library for lifelong Minecraft exploration, and LearnAct creates and revises callable actions for interactive environments \cite{wang2024voyager,zhao2024learnact}. SkillWeaver discovers, practices, and distills website-specific APIs for later web tasks \cite{zheng2025skillweaver}. These systems emphasize reliable invocation and reuse, but their artifacts are executable capabilities rather than fixed natural-language debugging procedures.

Procedural memory is now an explicit object of study. Memp compares fine-grained instructions with higher-level script-like abstractions and studies how a memory repository is built, retrieved, updated, corrected, and deprecated \cite{fang2026memp}. Skill-Pro represents skills through activation, execution, and termination conditions and maintains them through verification and score-based updates \cite{mi2026skillpro}. LEGOMem decomposes workflow trajectories into modular memory units allocated across orchestrators and task agents \cite{han2026legomem}. Our work therefore does not claim to introduce procedural memory. It studies a narrower artifact and evaluation setting: low-shot induction of a frozen natural-language \texttt{SKILL.md} for held-out debugging tasks.

\subsection{Language-agent debugging and benchmarks}
\label{sec:related_debugging}

Repository-level software engineering provides a demanding test of agent transfer. SWE-bench evaluates whether language models can resolve real GitHub issues under repository tests, while SWE-agent shows that the agent-computer interface materially affects repair behavior \cite{jimenez2024swebench,yang2024sweagent}. ReasoningBank connects this domain to memory research by evaluating reusable reasoning memories on SWE-bench Verified and reporting success together with interaction steps and token consumption \cite{ouyang2026reasoningbank}. It also finds that retrieving excessive or noisy experience can reduce performance, making it the closest prior comparison for efficiency and harmful memory in software repair.

Reproducible bug databases support more controlled debugging studies. BugsInPy packages real Python defects with environments and tests, and PyBugHive expands this line with a larger manually validated collection of reproducible Python bugs \cite{widyasari2020bugsinpy,antal2024pybughive}. These benchmarks provide bug provenance and executable validation, but they do not define how an agent should induce or evaluate transferable memory. We use PyBugHive to construct a within-family setting in which training trajectories and held-out failures remain related without being identical.

\subsection{Evaluation gap}
\label{sec:related_gap}

Prior work already evaluates more than final success. Cross-task memory systems report actions, steps, API calls, tokens, or task efficiency, and several papers recognize that stale, incompatible, or noisy memory can be harmful \cite{chen2024automanual,wang2025workflowmemory,fang2026memp,ouyang2026reasoningbank}. The remaining gap is therefore not the existence of procedural memory, efficiency measurement, or memory degradation in isolation.

Within our verified corpus, we did not identify a study that jointly evaluates five elements: (1) a fixed natural-language procedure induced from a few debugging trajectories, (2) transfer to held-out bugs in a controlled software family, (3) repair success together with diagnosis-edit-test cycles, token cost, and failed patches, (4) explicit per-task attribution of negative transfer to the memory condition, and (5) controls for generic advice, reflection-style memory, memory length, and procedural organization. Our contribution targets this joint evaluation gap. The empirical study remains deliberately narrow, so the claim concerns a reproducible protocol and an initial within-family hard-smoke evaluation rather than general superiority over prior memory systems.
